\documentclass[11pt,a4paper]{article}

% ====== Paquetes ======
\usepackage[utf8]{inputenc}
\usepackage[T1]{fontenc}
\usepackage[spanish]{babel}
\usepackage{geometry}
\geometry{margin=2.4cm}
\usepackage{listings}
\usepackage{xcolor}
\usepackage{hyperref}
\usepackage{enumitem}
\usepackage{titlesec}
\usepackage{amsmath}
\usepackage{tcolorbox}

% ====== Hyperref ======
\hypersetup{
    colorlinks=true,
    linkcolor=blue!60!black,
    urlcolor=blue!60!black,
    citecolor=blue!60!black
}

% ====== Estilo de código Python ======
\definecolor{codegray}{rgb}{0.95,0.95,0.95}
\definecolor{codekey}{rgb}{0.0,0.0,0.6}
\definecolor{codestr}{rgb}{0.58,0.0,0.0}
\definecolor{codecomment}{rgb}{0.13,0.55,0.13}

\lstdefinestyle{pythonstyle}{
    language=Python,
    backgroundcolor=\color{codegray},
    basicstyle=\ttfamily\footnotesize,
    keywordstyle=\color{codekey}\bfseries,
    commentstyle=\color{codecomment}\itshape,
    stringstyle=\color{codestr},
    showstringspaces=false,
    frame=single,
    rulecolor=\color{gray!40},
    breaklines=true,
    tabsize=4,
    numbers=left,
    numberstyle=\tiny\color{gray},
    numbersep=8pt,
    captionpos=b
}
\lstset{style=pythonstyle}

% ====== Caja de notas pedagógicas ======
\newtcolorbox{notabox}[1][]{
    colback=yellow!8,
    colframe=orange!60!black,
    boxrule=0.5pt,
    arc=2pt,
    left=6pt,right=6pt,top=4pt,bottom=4pt,
    title=Nota de aprendizaje,
    fonttitle=\bfseries\small,
    #1
}

% ====== Título ======
\title{\textbf{Plan Maestro de Desarrollo} \\[0.3em] \Large Proyecto Hill Climb AI \\[0.2em] \normalsize Versión 2.0 --- Revisada y comentada}
\author{}
\date{}

\begin{document}
\maketitle
\vspace{-2em}
\tableofcontents
\newpage

% ============================================================
\section{Objetivo General}
% ============================================================

Desarrollar un entorno inspirado en \textit{Hill Climb Racing} utilizando únicamente Python. El proyecto tiene dos componentes principales:

\begin{enumerate}
    \item \textbf{Simulación jugable} con física basada en restricciones rígidas.
    \item \textbf{Entrenamiento de una IA} mediante neuroevolución para aprender a conducir automáticamente.
\end{enumerate}

La simulación se ejecuta visualmente en tiempo real y muestra el aprendizaje de la IA mientras evoluciona generación tras generación.

\textbf{El objetivo final NO es construir un videojuego comercial}, sino un entorno de aprendizaje por refuerzo (vía neuroevolución) modular, extensible y presentable como proyecto de portafolio en ciencia de datos.

\begin{notabox}
La neuroevolución es una alternativa al aprendizaje por refuerzo clásico (Q-learning, PPO). En lugar de calcular gradientes para mejorar la política, se mantiene una \textit{población} de redes neuronales, se evalúa cuál juega mejor, y los mejores se ``reproducen'' mutando sus pesos. Es ideal para problemas donde la recompensa es escasa o ruidosa, y se puede paralelizar bien.
\end{notabox}

% ============================================================
\section{Tecnologías del Proyecto}
% ============================================================

Instalación inicial:

\begin{lstlisting}[language=bash]
pip install pygame pymunk numpy torch matplotlib
\end{lstlisting}

\noindent\textbf{Función de cada librería:}

\begin{itemize}[leftmargin=*]
    \item \textbf{pygame}: renderizado de gráficos 2D, captura de eventos de teclado/ratón, UI.
    \item \textbf{pymunk}: motor de física 2D (basado en Chipmunk). Maneja cuerpos rígidos, articulaciones, colisiones y gravedad.
    \item \textbf{numpy}: operaciones vectoriales rápidas (lookahead de terreno, normalización de estados).
    \item \textbf{torch}: implementación de la red neuronal. Aunque no usaremos backpropagation, PyTorch facilita el manejo de tensores y será útil si extiendes el proyecto a DRL más adelante.
    \item \textbf{matplotlib}: gráficas de evolución del fitness, distribuciones de mutación, métricas de entrenamiento.
\end{itemize}

\begin{notabox}
\textbf{Sobre PyTorch en neuroevolución:} como NO usamos gradientes, todos los tensores de la red deben crearse con \texttt{requires\_grad=False}, y los forward passes deben envolverse en \texttt{with torch.no\_grad():}. Esto evita que PyTorch construya el grafo computacional y mejora el rendimiento. Es una excelente oportunidad para entender cómo PyTorch separa el cálculo del autograd.
\end{notabox}

% ============================================================
\section{Arquitectura Final}
% ============================================================

La estructura NO debe modificarse durante el proyecto.

\begin{lstlisting}[language=bash]
hill_climb_ai/
    main.py
    config/
        settings.py
    game/
        environment.py
        terrain.py
        vehicle.py
        player.py
        coin.py
        checkpoint.py
        physics.py
        ui.py
        camera.py
    ai/
        neural_network.py
        genome.py
        genetic_algorithm.py
        trainer.py
        reward_system.py
    utils/
        helpers.py
        logger.py
        math_utils.py
    assets/
        images/
        fonts/
    data/
        saved_models/
        statistics/
    experiments/
        physics_tests.ipynb
        reward_tests.ipynb
        plots.ipynb
\end{lstlisting}

% ============================================================
\section{Responsabilidad de Cada Archivo}
% ============================================================

\subsection{main.py}

\textbf{Responsabilidades:}
\begin{itemize}
    \item Inicializar \texttt{pygame} y crear ventana.
    \item Parsear argumentos de línea de comandos (\texttt{--mode play|train|watch}, \texttt{--render true|false}).
    \item Ejecutar el game loop principal.
    \item Controlar FPS según el modo (60 en juego/observación, sin límite en entrenamiento headless).
    \item Lanzar el entrenamiento llamando a \texttt{ai/trainer.py}.
\end{itemize}

\textbf{Prohibido:} crear lógica física, calcular recompensas, programar IA.

\subsection{config/settings.py}

Contiene únicamente constantes globales. Ningún import de otros módulos del proyecto.

\begin{lstlisting}
# === Ventana y rendering ===
SCREEN_WIDTH = 1280
SCREEN_HEIGHT = 720
FPS = 60

# === Fisica ===
GRAVITY = 900           # px/s^2
WHEEL_FRICTION = 1.5
CHASSIS_MASS = 5.0
WHEEL_MASS = 1.0

# === Episodio ===
MAX_TIME = 20.0         # segundos sin checkpoint -> muerte
CHECKPOINT_TIME = 10.0  # segundos que anade un checkpoint
TARGET_DISTANCE = 100   # px minimos para considerar avance

# === Red neuronal ===
NN_INPUTS = 14
NN_HIDDEN = [16, 12]
NN_OUTPUTS = 2
LOOKAHEAD_DISTANCES = [30, 80, 150, 250]  # px frente al vehiculo

# === Algoritmo genetico ===
POPULATION_SIZE = 50
ELITISM_COUNT = 3        # mejores individuos que pasan sin mutar
MUTATION_RATE = 0.1      # probabilidad por peso
MUTATION_SIGMA = 0.2     # desviacion estandar del ruido gaussiano
TOURNAMENT_SIZE = 3
CROSSOVER_RATE = 0.7

# === Persistencia ===
SAVE_EVERY_N_GEN = 5
MODEL_DIR = "data/saved_models"
STATS_DIR = "data/statistics"
\end{lstlisting}

\begin{notabox}
Todas las magnitudes con unidad están comentadas. Esto es importante en física simulada: confundir píxeles/segundo con metros/segundo es la causa \#1 de comportamientos raros en pymunk.
\end{notabox}

\subsection{game/environment.py}

Archivo central del juego. Orquesta todos los componentes.

\textbf{Responsabilidades:}
\begin{itemize}
    \item Inicializar terreno, vehículo, jugador, monedas, checkpoints.
    \item Llamar a \texttt{physics.py} para avanzar la simulación un \texttt{dt}.
    \item Verificar condición de muerte (jugador toca el suelo o tiempo expirado).
    \item Actualizar score y gestionar colecciones (monedas, checkpoints alcanzados).
    \item Exponer \texttt{get\_state()} que devuelve el vector de 14 entradas para la IA.
    \item Exponer \texttt{step(action)} (interfaz tipo Gym): recibe acción, avanza un frame, retorna \texttt{(state, reward, done, info)}.
\end{itemize}

\subsection{game/terrain.py}

Generación del mapa.

\textbf{Debe:}
\begin{itemize}
    \item Crear colinas y pendientes manteniendo continuidad.
    \item Aceptar una \textbf{semilla} (seed) para reproducibilidad.
    \item Exponer \texttt{height\_at(x)} para consultar altura del terreno en cualquier coordenada x (necesario para el lookahead de la red).
    \item Exponer \texttt{slope\_at(x)} para consultar pendiente local.
\end{itemize}

\textbf{Primera versión:} terreno fijo (lista de puntos predefinida).\\
\textbf{Versión posterior:} terreno procedural por funciones sinusoidales superpuestas (Perlin noise opcional).

\begin{notabox}
La semilla del terreno debe ser \textbf{idéntica para todos los individuos de una misma generación}. Si cada uno enfrenta un terreno distinto, el fitness no es comparable. Se cambia la semilla solo al pasar a la siguiente generación (o cada N generaciones).
\end{notabox}

\subsection{game/vehicle.py}

Implementación física del vehículo en pymunk.

\textbf{Debe contener:}
\begin{itemize}
    \item Chasis (cuerpo rígido rectangular).
    \item Dos ruedas (cuerpos circulares unidos al chasis por articulaciones tipo \texttt{pin joint} o \texttt{groove joint}).
    \item Motor en cada rueda con torque controlable.
    \item Tracking de inclinación y velocidad angular del chasis.
\end{itemize}

\textbf{Interfaz mínima:}

\begin{lstlisting}
def accelerate(intensity: float) -> None: ...   # intensity in [0, 1]
def brake(intensity: float) -> None: ...        # intensity in [0, 1]
def update(dt: float) -> None: ...
def get_telemetry() -> dict: ...                # vx, vy, angle, omega, contacts
\end{lstlisting}

\subsection{game/player.py}

Representa al conductor encima del chasis.

\textbf{Regla principal:} si el cuerpo del conductor toca el suelo, el episodio termina (\texttt{done = True}).

Implementado como un cuerpo rígido pequeño unido al chasis por un \texttt{pivot joint} rígido. La colisión con el terreno se detecta usando \texttt{collision\_handlers} de pymunk.

\subsection{game/coin.py}

Cada moneda contiene:
\begin{itemize}
    \item \texttt{position}: (x, y) en coordenadas del mundo.
    \item \texttt{value}: puntos otorgados al recogerla (por defecto 1).
    \item \texttt{collected}: booleano.
\end{itemize}

Las monedas se distribuyen por el terreno con espaciado pseudo-aleatorio.

\subsection{game/checkpoint.py}

Cada checkpoint:
\begin{itemize}
    \item Agrega \texttt{CHECKPOINT\_TIME} segundos al contador del episodio.
    \item Otorga una recompensa fija (\texttt{+200}).
    \item Reinicia el contador de inactividad.
\end{itemize}

\subsection{game/physics.py}

\textbf{Antes ambiguo, ahora con rol definido.}

\textbf{Responsabilidades:}
\begin{itemize}
    \item Crear y mantener el \texttt{pymunk.Space} global.
    \item Configurar gravedad, damping y constantes de simulación desde \texttt{settings.py}.
    \item Exponer \texttt{step(dt)} que avanza la simulación.
    \item Registrar \texttt{collision\_handlers} (vehículo-suelo, jugador-suelo, vehículo-moneda).
\end{itemize}

\textbf{Prohibido:} dibujar nada, contener lógica de juego.

\subsection{game/ui.py}

Renderiza información sobre la pantalla.

\textbf{Mostrar:}
\begin{itemize}
    \item Score actual, generación, distancia, tiempo restante, mejor fitness histórico.
    \item Mini-gráfica de evolución del fitness (opcional, post fase 6).
\end{itemize}

\textbf{También:} botones de acelerar/frenar (modo manual).

\subsection{game/camera.py}

Sigue al vehículo horizontalmente con un pequeño \textit{lerp} para suavizar.

\textbf{NO debe contener lógica de juego.} Solo transformaciones de coordenadas mundo $\rightarrow$ pantalla.

\subsection{ai/neural\_network.py}

Red neuronal feedforward implementada en PyTorch.

\textbf{Arquitectura:}

\begin{center}
14 entradas $\rightarrow$ 16 (tanh) $\rightarrow$ 12 (tanh) $\rightarrow$ 2 (sigmoide)
\end{center}

\textbf{Entradas (vector de 14 dimensiones):}

\begin{enumerate}
    \item Velocidad horizontal $v_x$ normalizada.
    \item Velocidad vertical $v_y$ normalizada.
    \item Ángulo del chasis $\theta$ (radianes, en $[-\pi, \pi]$).
    \item Velocidad angular $\omega$.
    \item Contacto rueda izquierda (0 o 1).
    \item Contacto rueda derecha (0 o 1).
    \item Altura terreno relativa a +30 px frente al vehículo.
    \item Altura terreno relativa a +80 px.
    \item Altura terreno relativa a +150 px.
    \item Altura terreno relativa a +250 px.
    \item Pendiente del terreno en la posición actual.
    \item $\Delta x$ a moneda más cercana (normalizado).
    \item $\Delta y$ a moneda más cercana (normalizado).
    \item Tiempo restante normalizado (0 a 1).
\end{enumerate}

\textbf{Salidas:}
\begin{enumerate}
    \item Intensidad de aceleración (sigmoide, $[0, 1]$).
    \item Intensidad de frenado (sigmoide, $[0, 1]$).
\end{enumerate}

\textbf{Importante:} la red debe crearse con \texttt{requires\_grad=False} en todos sus parámetros, y los forward passes envolverse en \texttt{torch.no\_grad()}.

\begin{lstlisting}
import torch
import torch.nn as nn

class PolicyNet(nn.Module):
    def __init__(self, n_in=14, hidden=[16, 12], n_out=2):
        super().__init__()
        layers = []
        prev = n_in
        for h in hidden:
            layers += [nn.Linear(prev, h), nn.Tanh()]
            prev = h
        layers += [nn.Linear(prev, n_out), nn.Sigmoid()]
        self.net = nn.Sequential(*layers)
        # Desactivar gradientes (no usamos backprop)
        for p in self.parameters():
            p.requires_grad = False

    def forward(self, x):
        with torch.no_grad():
            return self.net(x)
\end{lstlisting}

\subsection{ai/genome.py}

Un \texttt{Genome} es un individuo de la población. Contiene:
\begin{itemize}
    \item Su red neuronal (\texttt{PolicyNet}).
    \item Su fitness (inicialmente 0).
    \item Métodos: \texttt{get\_weights()}, \texttt{set\_weights(vector)}, \texttt{mutate()}, \texttt{copy()}.
\end{itemize}

Los pesos se manejan como un \textbf{vector plano (1D)} de NumPy para facilitar crossover y mutación. La conversión vector $\leftrightarrow$ red se hace con \texttt{torch.nn.utils.parameters\_to\_vector} y \texttt{vector\_to\_parameters}.

\subsection{ai/genetic\_algorithm.py}

Implementa los operadores evolutivos.

\textbf{Selección:} torneo de tamaño \texttt{TOURNAMENT\_SIZE}.

\textbf{Crossover:} uniforme, gen por gen (probabilidad 50\% de tomar el peso del padre A o B).

\textbf{Mutación:} para cada peso, con probabilidad \texttt{MUTATION\_RATE}, sumar ruido gaussiano $\mathcal{N}(0, \sigma)$ con $\sigma = $ \texttt{MUTATION\_SIGMA}.

\textbf{Elitismo:} los \texttt{ELITISM\_COUNT} mejores individuos pasan sin modificación a la siguiente generación. Esto evita perder al mejor genoma por mala suerte en crossover.

\textbf{Nueva población:} elite + descendencia hasta completar \texttt{POPULATION\_SIZE}.

\subsection{ai/trainer.py}

Orquestador del entrenamiento.

\begin{lstlisting}
Crear poblacion inicial (o cargar desde disco)
loop generaciones:
    para cada genoma:
        ejecutar episodio en environment
        calcular fitness final
    ordenar poblacion por fitness
    guardar metricas
    si generacion % SAVE_EVERY_N_GEN == 0:
        guardar poblacion
    aplicar operadores geneticos para crear siguiente generacion
\end{lstlisting}

\subsection{ai/reward\_system.py}

Define la función de fitness del episodio.

\textbf{Fitness final del episodio:}

\[
\text{fitness} = d_{\max} + 50 \cdot n_{\text{monedas}}
\]

donde $d_{\max}$ es la distancia horizontal máxima alcanzada y $n_{\text{monedas}}$ es la cantidad de monedas recogidas.

\textbf{Recompensa instantánea (para debugging y análisis, no para el GA):}

\begin{lstlisting}
reward_step = 0
reward_step += delta_distance * 0.1
reward_step += 50  if coin_collected else 0
reward_step += 200 if checkpoint_reached else 0
reward_step -= 200 if player_touched_ground else 0
\end{lstlisting}

\begin{notabox}
\textbf{Distinción clave:} las recompensas instantáneas (\texttt{reward\_step}) son útiles si después extiendes el proyecto a Q-learning o PPO. Para el algoritmo genético solo importa el \textbf{fitness final del episodio}, que es un único escalar por individuo.
\end{notabox}

\subsection{utils/helpers.py}

Funciones de propósito general: conversión de coordenadas, normalización de vectores, clamp, lerp.

\subsection{utils/logger.py}

Logger del entrenamiento.

\textbf{Responsabilidades:}
\begin{itemize}
    \item Escribir un \texttt{.csv} por generación con: generación, fitness máx, fitness medio, fitness mín, desviación estándar, mejor distancia, mejor cantidad de monedas.
    \item Imprimir en consola un resumen formateado por generación.
\end{itemize}

\subsection{utils/math\_utils.py}

Funciones matemáticas reusables: cálculo de pendiente discreta, distancia euclidiana, normalización min-max, clipping de ángulos a $[-\pi, \pi]$.

% ============================================================
\section{Persistencia}
% ============================================================

El entrenamiento NO puede perderse. El sistema debe ser robusto a cierres inesperados.

\textbf{Qué guardar:}
\begin{itemize}
    \item Vector de pesos del mejor individuo (\texttt{best\_genome.pt}).
    \item Población completa en formato pickle (\texttt{population\_gen\_N.pkl}).
    \item Métricas históricas (\texttt{stats.csv}).
    \item Número de generación actual.
\end{itemize}

\textbf{Estructura de carpetas:}

\begin{lstlisting}[language=bash]
data/
    saved_models/
        best_genome.pt
        population_gen_5.pkl
        population_gen_10.pkl
        ...
    statistics/
        stats.csv
\end{lstlisting}

\textbf{Política de guardado:}

\begin{lstlisting}
if generation % SAVE_EVERY_N_GEN == 0:
    save_population(generation)
    save_best_genome()
save_stats_row(generation, metrics)  # cada generacion
\end{lstlisting}

\textbf{Al iniciar:}

\begin{lstlisting}
if exists(latest_population_file):
    population, start_gen = load_population()
else:
    population = create_random_population()
    start_gen = 0
\end{lstlisting}

% ============================================================
\section{Plan de Desarrollo}
% ============================================================

\subsection{FASE 1: Preparación}

\textbf{Objetivo:} crear la estructura completa del proyecto.

\textbf{Humano:}
\begin{enumerate}
    \item Crear repositorio y carpeta raíz.
    \item Crear todas las subcarpetas y archivos vacíos según la arquitectura.
    \item Crear entorno virtual e instalar librerías.
    \item Inicializar git y hacer primer commit.
\end{enumerate}

\textbf{IA:} aún no programa. Solo verifica que la arquitectura esté correctamente creada.

\textbf{Resultado esperado:} proyecto vacío pero perfectamente organizado.

\subsection{FASE 2: Juego Manual}

\textbf{Objetivo:} construir la simulación sin IA, jugable con teclado.

\textbf{Orden estricto de implementación:}
\begin{enumerate}
    \item \texttt{config/settings.py} (constantes).
    \item \texttt{game/physics.py} (Space de pymunk).
    \item \texttt{game/terrain.py} (terreno fijo primero).
    \item \texttt{game/vehicle.py} (chasis + ruedas + torque).
    \item \texttt{game/player.py} (conductor + colisión con suelo).
    \item \texttt{game/camera.py} (seguimiento horizontal).
    \item \texttt{game/ui.py} (HUD básico).
    \item \texttt{game/environment.py} (integración de todo).
    \item \texttt{main.py} (game loop).
\end{enumerate}

\textbf{Regla obligatoria:} probar cada módulo antes de continuar al siguiente.

\textbf{Resultado esperado:} se puede jugar manualmente con teclado (flechas izq/der o A/D), el vehículo se mueve sobre el terreno, la cámara sigue, y si el conductor toca el suelo el episodio termina.

\subsection{FASE 3: Sistema de Recompensas}

Añadir:
\begin{itemize}
    \item \texttt{game/coin.py} con distribución de monedas.
    \item \texttt{game/checkpoint.py} con extensión de tiempo.
    \item Score y tiempo en \texttt{ui.py}.
    \item Esqueleto de \texttt{ai/reward\_system.py} (aún sin IA).
\end{itemize}

\textbf{Humano:} verificar manualmente que recoger monedas suma score, que pasar por checkpoints extiende tiempo, y que el episodio termina correctamente al agotar tiempo o tocar suelo.

\textbf{Resultado esperado:} juego completo y funcional sin IA.

\subsection{FASE 4: Estados de IA}

Implementar en \texttt{environment.py} la función \texttt{get\_state()} que retorne el vector de 14 entradas.

\textbf{Humano:} ejecutar el juego manualmente y graficar (en \texttt{experiments/plots.ipynb}) cómo varían las 14 variables a lo largo de un episodio. Verificar:
\begin{itemize}
    \item Que los rangos sean razonables (idealmente $[-1, 1]$ tras normalizar).
    \item Que las variables de lookahead respondan correctamente al terreno.
    \item Que los contactos de ruedas sean estables (no parpadeen).
\end{itemize}

\subsection{FASE 5: Red Neuronal}

Implementar \texttt{ai/neural\_network.py} y \texttt{ai/genome.py}.

\textbf{Prueba clave:} crear una red con pesos aleatorios, conectarla al juego, y verificar que el vehículo se mueva (caóticamente, pero se mueva).

\textbf{Resultado esperado:} la IA controla el vehículo aunque sin aprender aún.

\subsection{FASE 6: Algoritmo Genético}

Implementar \texttt{ai/genetic\_algorithm.py} y \texttt{ai/trainer.py}.

\textbf{Pruebas obligatorias:}
\begin{itemize}
    \item Ejecutar 5 generaciones cortas (población pequeña) y verificar que el fitness promedio aumente.
    \item Graficar curva de evolución (mejor, promedio, peor por generación).
    \item Verificar elitismo: el mejor individuo nunca empeora entre generaciones.
\end{itemize}

\textbf{Resultado esperado:} aprendizaje visible. Tras $\sim$30--50 generaciones la IA debe avanzar consistentemente.

\subsection{FASE 7: Persistencia}

Agregar guardado y carga en \texttt{trainer.py}.

\textbf{Prueba de robustez:}
\begin{enumerate}
    \item Entrenar 20 generaciones.
    \item Cerrar el programa.
    \item Reiniciar: el entrenamiento debe continuar desde la generación 20.
    \item Cargar el mejor genoma en modo \texttt{watch} y observar su desempeño.
\end{enumerate}

\subsection{FASE 8: Optimización y Extensiones (Opcional)}

Funcionalidades opcionales para enriquecer el proyecto:

\begin{itemize}
    \item \textbf{Modo headless de entrenamiento.} Flag \texttt{--render=False} que desactive \texttt{pygame.display.update()}. Permite correr cientos de generaciones rápido y solo renderizar cuando se quiere observar.
    \item \textbf{Evaluación en paralelo.} Múltiples vehículos en el mismo \texttt{pymunk.Space} con filtros de colisión (no chocan entre sí). Acelera enormemente la evaluación de la población.
    \item \textbf{Terreno procedural.} Suma de sinusoides o Perlin noise con semilla controlada.
    \item \textbf{NEAT (NeuroEvolution of Augmenting Topologies).} Evoluciona la topología además de los pesos. Más potente pero más complejo. Existe la librería \texttt{neat-python}, o se puede implementar desde cero como ejercicio.
    \item \textbf{Obstáculos dinámicos} (rocas, rampas, gaps).
    \item \textbf{Métricas avanzadas:} diversidad genética (varianza de pesos en la población), tasa de convergencia, análisis de neuronas más activas.
    \item \textbf{Guardado automático del mejor video} (grabar el episodio del mejor genoma de cada generación con \texttt{imageio}).
\end{itemize}

% ============================================================
\section{Reglas de Vibe Coding}
% ============================================================

Estas reglas son obligatorias y aplican a cada interacción con Claude Code.

\textbf{Nunca pedir:}
\begin{itemize}
    \item ``Hazme todo el proyecto.''
    \item ``Implementa varias fases a la vez.''
\end{itemize}

\textbf{Siempre pedir:}
\begin{itemize}
    \item ``Explícame cómo funciona X antes de implementarlo.''
    \item ``Implementa solamente \texttt{terrain.py} usando la interfaz ya existente.''
    \item ``Muéstrame un ejemplo mínimo de Y antes de integrarlo.''
\end{itemize}

\textbf{Reglas obligatorias:}
\begin{enumerate}
    \item \textbf{Un archivo por iteración.} Nunca dos a la vez.
    \item \textbf{Integrar y probar antes de seguir.} Cada módulo nuevo se ejecuta y se verifica visualmente o con prints/asserts.
    \item \textbf{Commit por módulo.} Cada archivo terminado y probado es un commit.
    \item \textbf{Preguntar el ``por qué'' antes del ``cómo''.} Antes de copiar código generado, pedir que se explique la decisión de diseño.
    \item \textbf{Comentarios pedagógicos.} Pedir a la IA que comente el código a nivel de aprendizaje, no solo a nivel de documentación.
\end{enumerate}

% ============================================================
\section{Resultado Esperado}
% ============================================================

Al finalizar el proyecto se obtiene:

\begin{itemize}
    \item Un juego tipo Hill Climb funcional, jugable manualmente.
    \item Una IA entrenada por neuroevolución capaz de avanzar consistentemente y recoger monedas.
    \item Visualización en tiempo real del aprendizaje generación tras generación.
    \item Guardado persistente del entrenamiento, robusto a reinicios.
    \item Arquitectura modular y limpia, fácil de extender.
    \item Notebooks de experimentos con gráficas de evolución del fitness.
    \item Un proyecto sólido de portafolio en ciencia de datos / RL aplicado.
\end{itemize}

\vspace{2em}
\hrule
\vspace{0.5em}
\noindent\textit{Documento revisado. Versión 2.0. Listo para la fase de implementación.}

\end{document}
